\chapter{Theoretical Background}\label{ch:theoretical-background}



\section{State of the Art}


\subsection{Evolution of Seizure Prediction and Algorithmic Development}
The pursuit of identifying seizure precursors (preictal states)  from the electroencephalogram began in the 1970s with \textit{linear approaches}, such as autoregressive modeling and power-spectrum analysis of spike-wave bursts \parencite{mormann_seizure_2007}.
Early research also focused on the predictive value of \textit{interictal spike rates}, though results were often contradictory; while some early studies suggested focal spiking decreased before seizures, subsequent investigations on more extended databases showed no systematic or reproducible changes \parencite{mormann_seizure_2007}.

In the 1990s, the field shifted toward \textit{nonlinear dynamics}, treating the brain as a chaotic system \parencite{mormann_seizure_2007, kuhlmann_epilepsyecosystemorg_2018}.
Researchers utilized measures like the largest Lyapunov exponent to identify a decrease in chaoticity in the minutes preceding a seizure, alongside the dynamical similarity index and correlation density to characterize a transitional preictal phase \parencite{mormann_seizure_2007}.
This era was marked by significant optimism, with many studies reporting high sensitivity on small, highly selected datasets \parencite{mormann_seizure_2007, muller_coherent_2022}.
However, a period of \textit{skepticism} followed in the early 2000s as more rigorous evaluations on unselected, continuous multi-day recordings failed to reproduce these optimistic findings, highlighting issues with overfitting, lack of specificity evaluation, and the absence of out-of-sample statistical validation \parencite{mormann_seizure_2007, freestone_forward-looking_2017}.

To overcome these reproducibility challenges and address the \enquote{ceiling effect} in predictive performance, the field turned to \textit{crowdsourcing and open competitions} using standardized, high-quality datasets obtained using the NeuroVista device, the first implantable seizure prediction system \parencite{cook_prediction_2013}. Platforms like \textit{Kaggle} hosted the American Epilepsy Society Seizure Prediction Challenge, providing \gls{ieeg} recordings from both humans and canines \parencite{brinkmann_crowdsourcing_2016}.
Subsequent efforts by \textit{epilepsyecosystem.org} continued this trend, enabling researchers to compare algorithms on ultra-long-term human recordings \parencite{kuhlmann_epilepsyecosystemorg_2018}.
These competitions demonstrated that \textit{machine learning} could achieve better-than-chance performance \parencite{brinkmann_crowdsourcing_2016}.
Winning solutions typically utilized a variety of \enquote{hand-crafted} features---including band power spectra, statistical moments (skewness, kurtosis), and cross-channel linear correlations in both time and frequency domains---classified by \textit{Support Vector Machines (SVM)}, \textit{Random Forests}, or \textit{k-nearest neighbors} \parencite{kuhlmann_epilepsyecosystemorg_2018, muller_coherent_2022}.

More recently, \textit{deep learning} has emerged as a dominant paradigm, with \textit{\gls{cnn}} and \textit{\gls{lstm}} networks used to extract features directly from raw data, bypassing the need for manual feature engineering \parencite{kuhlmann_epilepsyecosystemorg_2018, muller_coherent_2022}.
Looking forward, \textcite{freestone_forward-looking_2017} suggest that the field is shifting from a deterministic preictal process underlying the seizure prediction paradigm toward a probabilistic proictal state, thereby accepting the inherent uncertainty and moving toward seizure forecasting. 


\subsection{Seizure forecasting}
The shift toward seizure forecasting was initiated following several seminal studies that identified a relationship between circadian and multiday cycles and the occurrence of epileptic seizures. Seizures were found to preferentially occur at specific phases of these cycles, thereby establishing prolonged periods of elevated seizure propensity—the proictal state—rather than direct seizure precursors, i.e., the preictal state.


\subsection{Precursors and Mechanisms of Ictogenesis}
Understanding the mechanisms by which a \enquote{normal} brain transitions into a seizure---a process known as \textit{ictogenesis}---is fundamental to predictive modeling \parencite{mormann_seizure_2007}.
While some patients report clinical prodromi such as behavioral or mood changes, the field operates on the assumption that these transitions are best tracked by \textit{\gls{eeg}}, which can capture measurable shifts in brain dynamics hours or even days before a clinical event in many patients \parencite{freestone_forward-looking_2017, cook_prediction_2013}.

% emphasize the discrepancy of hypotheses in the literature and mention other hypothesis (from Mormann 2007)
On a network level, focal seizures are hypothesized to arise from \textit{\enquote{burster} neurons} that abnormally recruit and entrain neighbors into a \enquote{critical mass}.
This build-up is viewed not merely as an increasing count of neurons but as an unfolding process of increasing \textit{critical interactions} within a pathological functional network \parencite{mormann_seizure_2007}.

% rewrite as Another group has tried to ... (there's not necessarily a direct connection)
A key indicator of this transition is \enquote{\textit{critical slowing}}, which refers to the tendency of a system's regulatory mechanisms to slow down, thereby making the system take longer to return to equilibrium after perturbations.
This is viewed as a universal signature of a system losing its resilience as it approaches a tipping point or state transition.
In the context of epilepsy, this is often characterized by an increase in the autocorrelation of the \gls{eeg} or a lengthening of electrical evoked potentials \parencite{freestone_forward-looking_2017, maturana_critical_2020}.

This transition is frequently framed through the lens of \textit{bi/multistability}, where a stable healthy attractor (e.g., wakefulness or sleep) coexists with a pathological seizure state.
Seizures may occur due to deterministic variations in underlying parameters or random \enquote{noisy} fluctuations in background activity that drive the brain's state across the energy barrier separating these attractors.
Because the brain possesses internal \textit{homeostatic regulatory mechanisms}, it is possible for the brain to enter a state of \textit{high susceptibility} without a clinical event actually occurring, as these mechanisms may correct the abnormal activity \parencite{freestone_forward-looking_2017}.

This conceptualization leads to the idea of a \textit{proictal state}---a period of heightened seizure likelihood that does not always terminate in an ictal event \parencite{muller_coherent_2022}.
Furthermore, research into biomarkers like \textit{interictal spikes} and \textit{\glspl{hfo}} has revealed that while these indicators have clear implications for understanding ictogenesis, their dynamics are highly patient-specific \parencite{freestone_forward-looking_2017}.

Consequently, there is an increasing recognition that \enquote{one-size-fits-all} algorithms are insufficient, necessitating long-term monitoring to account for individual variability in seizure-generating mechanisms and the diverse populations of seizures that can coexist within a single patient \parencite{freestone_forward-looking_2017, yang_seizure_2024}.



\section{Electroencephalography (EEG)}\label{sec:eeg-background}

\textbf{\Gls{eeg}} is an electrophysiological technique used to record the electrical activity arising from brains.
It primarily measures the extracellular ion flow generated by similarly oriented groups of cerebral cortical neurons---especially pyramidal neurons---near the scalp where the recording electrodes are placed.
The measurement represents the \textbf{summation of the inhibitory and excitatory postsynaptic potentials} (rather than action potentials, which are too fast to measure) of neurons firing synchronously.
Because these neurons are oriented perpendicularly to the brain's surface, their combined activity creates a dipole that can be detected by recording electrodes.
\gls{eeg} is characterized by its \textbf{exquisite temporal sensitivity}, providing a near real-time display of ongoing cerebral activity \parencite{st_louis_electroencephalography_2016}.

\subsection{Clinical Utility and Applications}
The primary utility of \gls{eeg} lies in the evaluation of \textbf{dynamic cerebral functioning}. 
It is widely used for several clinical indications \parencite{st_louis_electroencephalography_2016}:
\begin{itemize}
    \item \textbf{Epilepsy:} \gls{eeg} is the gold standard for diagnosing and confirming seizures and epilepsy. 
    It allows clinicians to identify characteristic ictal (during-seizure) alterations and \glspl{ied}, such as spikes and sharp waves, which help classify seizure types and epilepsy syndromes.
    \item \textbf{General monitoring:} It is used to monitor the \textit{depth of anesthesia} during surgery and to detect sudden changes in neural functioning that may indicate \textit{ischemia or infarction}.
    \item \textbf{Research and specialized diagnosis:} \gls{eeg} waveforms can be averaged to study \textit{\glspl{ep}} and \textit{\glspl{erp}}, which represent neural activity related to specific stimuli, assisting in the analysis of cognitive and sensory functioning.
\end{itemize}

\subsection{Frequency Bands}
\gls{eeg} activity is conventionally categorized into frequency\footnote{The exact frequency boundaries of each band vary in the literature and are based on \textcite{mormann_seizure_2007} here.} bandwidths, or bands, which vary according to the patient's state (e.g., wakefulness, sleep) and developmental stage \parencite{st_louis_electroencephalography_2016, mormann_seizure_2007, nayak_normal_2026}:

{
\sisetup{range-phrase = {--}, range-units = single}
\begin{itemize}
    \item \textbf{Delta $\delta$} (\SIrange{0.5}{4}{\hertz}): Predominant during deep sleep and considered pathological if present in awake adults.
    \item \textbf{Theta $\theta$} (\SIrange{4}{8}{\hertz}): Commonly seen during drowsiness and in early sleep stages.
    \item \textbf{Alpha $\alpha$} (\SIrange{8}{12}{\hertz}): Present during relaxed, awake states, especially in the occipital head region. Strengthens when eyes are closed and attenuates during intense activity or concentration.
    \item \textbf{Beta $\beta$} (\SIrange{12}{35}{\hertz}): Associated with alert wakefulness and often enhanced by most sedative medications.
    \item \textbf{Gamma $\gamma$} (\SIrange{35}{100}{\hertz}): Associated with active cognition. Is not commonly seen in scalp \gls{eeg} but is readily detected in intracranial recordings.
\end{itemize}
}


\subsection{EEG Recording Modalities}
The choice of electrode placement significantly impacts the recorded signal quality and spatial resolution \parencite{st_louis_electroencephalography_2016}:
\begin{itemize}
    \item \textbf{Surface (scalp) \gls{eeg}:} Non-invasive electrodes are placed on the skin using systems like the \textit{International 10--20 System} \parencite{klem_ten-twenty_1999}.
    However, the signal must traverse biological filters (brain, cerebrospinal fluid, meninges, skull, and skin), which reduce its amplitude and blur spatial information. 
    The skull acts as a \textit{spatial averager} and a dominant attenuator of high-frequency signals.
    \item \textbf{Video-\gls{eeg} Monitoring}: Considered the \enquote{gold standard} for characterizing paroxysmal spells and localizing seizures, this modality involves prolonged \gls{eeg} recording with simultaneous time-linked video. 
    Its primary value lies in the ability to correlate clinical behavioral changes (seizure semiology) with underlying electrographic activity. 
    It is essential for distinguishing true epileptic seizures from nonepileptic mimickers, such as psychogenic nonepileptic spells or syncope. 
    This monitoring is typically conducted in a specialized inpatient \gls{emu}, where antiepileptic drugs may be safely withdrawn to capture habitual seizures.
    \item \textbf{\Gls{ecog} / \gls{ieeg}:} Invasive monitoring involves placing electrodes directly on the cortex (subdural grids or strips) or within the brain (depth electrodes) \parencite{jayakar_diagnostic_2016}. 
    These provide much cleaner fidelity without muscle or movement artifacts, enabling precise localization of epileptic foci.
    \item \textbf{Subcutaneous (subscalp) \gls{eeg}:} This modality involves electrodes implanted under the skin but outside the skull, such as the \textit{UNEEG SubQ system} used in this study.
    Such systems offer a robust and unobtrusive alternative to traditional scalp \gls{eeg} \parencite{duun-henriksen_subdural_2013}.
\end{itemize}

Research comparing subcutaneous systems to standard surface electrodes shows that they achieve similar performance, as measured by \textit{\gls{psd}} and \textit{coherence analysis}.
Subcutaneous recordings often demonstrate \textit{higher signal quality} during rest, primarily because the electrodes provide \textit{higher signal amplitudes} and maintain a \textit{rigid connection} to the source, reducing artifacts caused by electrode-skin instability or movements \parencite{duun-henriksen_eeg_2015}.
While the skull still acts as a spatial averager for these signals, the subcutaneous approach avoids degradation of electrode--skin contact over time (common with surface electrodes) and is portable and inconspicuous, making it suitable for \enquote{ultra-long-term} monitoring over months.



\subsection{Seizure-Related Intervals in EEG Recordings}\label{subsec:background-periictal-intervals-in-eeg-recordings}

In the analysis of \gls{eeg} data for seizure prediction, the continuous recording is typically categorized into four distinct functional states or intervals based on their temporal relationship to a seizure event: ictal, pre-ictal, post-ictal, and inter-ictal \parencite{moghim_predicting_2014}:

\begin{itemize}
    \item \textbf{Ictal Interval}: This state represents the period of the actual seizure activity. It is characterized by sudden, abnormal, and synchronous neuronal discharges that typically manifest as high-amplitude, highly variable signal energy across EEG channels. While the duration varies, these events often last approximately three minutes.
    \item \textbf{Pre-ictal Interval}: This is the critical window immediately preceding a seizure (often defined as the 5 to 60 minutes prior to onset) where the brain undergoes a transition in dynamics.
    It has been shown that various clinical symptoms often occur before seizures, which might be usable as predictive markers.
    They include increased oxygen availability, cerebral blood flow, blood-oxygen-level-dependent signals, and changes in heart rate, and it is believed that an increasing number of critical interactions among neurons in the focal region unfold over time \parencite{moghim_predicting_2014}. 
    Signal properties such as accumulated energy and spectral edge frequency often show distinctive trends during this phase as the brain approaches a "critical mass" for ictogenesis.
    \item \textbf{Post-ictal Interval}: This period follows the seizure offset during which the brain recovers from the ictus. EEG signals during this state may show abnormal excitement or depressed activity as physiological functions return to baseline.
    \item \textbf{Inter-ictal Interval}: This state encompasses all non-seizure data that is temporally distant from both pre-ictal and post-ictal periods. It serves as the baseline for a patient's \enquote{normal} brain activity between events.
\end{itemize}

During the distinct intervals, electrophysiological changes in the brain affect the voltages recorded in the \gls{eeg} electrodes (\cref{fig:background-periictal-intervals}), thus, understanding the specific signal properties of these intervals is fundamental to developing predictive algorithms.
To distinguish these states, researchers analyze various quantitative descriptors, or features, extracted from the \gls{eeg} signal. 

\begin{figure}[hb]
    \centering
    \includegraphics[width=\linewidth]{ictal_intervals}
    \caption[Annotated Periictal Intervals of an \gls{ieeg} recording]{Annotated Ictal Intervals of an \gls{ieeg} recording. Source: \textcite{moghim_predicting_2014}}
    \label{fig:background-periictal-intervals}
\end{figure}
